\documentclass[UTF8,12pt]{ctexart}

\usepackage{amsmath,amssymb,amsthm}
\usepackage{geometry}
\usepackage{booktabs}
\usepackage{xcolor}
\usepackage{hyperref}

\geometry{a4paper, margin=2.5cm}

\title{乐观 UCB 与 Knowledge Gradient 的概率描述}
\author{}
\date{}

\begin{document}

\maketitle

\section{统一的不确定性建模}

考虑一组候选动作或选项 $a\in\mathcal A$。每个动作的真实价值为
\[
Q(a),
\]
但该价值未知。给定当前观测数据 $\mathcal D_t$，我们用后验分布描述其不确定性：
\[
Q(a)\mid \mathcal D_t \sim \mathcal N(\mu_a,\sigma_a^2),
\]
其中
\[
\mu_a
\]
表示当前对动作 $a$ 的价值均值估计，
\[
\sigma_a
\]
表示当前不确定性大小。

当前认知下的最优价值为
\[
V_t = \max_{a\in\mathcal A}\mu_a.
\]

\section{UCB：不确定性集合下的乐观最大化}

UCB 的核心思想是：

\[
\boxed{
\text{选择在乐观置信上界下可能最优的动作。}
}
\]

对每个动作 $a$，定义其上置信界：
\[
\mathrm{UCB}_t(a)
=
\mu_a+\beta_t\sigma_a,
\]
其中 $\beta_t>0$ 控制探索强度。

于是 UCB 策略为
\[
a_t^{\mathrm{UCB}}
=
\arg\max_{a\in\mathcal A}
\left(
\mu_a+\beta_t\sigma_a
\right).
\]

从概率角度看，如果后验近似为高斯分布，则
\[
\mathbb P
\left(
Q(a)\leq \mu_a+\beta_t\sigma_a
\right)
\]
可以被解释为某个较高置信水平。例如，当 $\beta_t=1.96$ 时，
\[
\mathbb P
\left(
Q(a)\leq \mu_a+1.96\sigma_a
\right)
\approx 0.975.
\]

因此 UCB 可以理解为：在高概率置信范围内，选择那个可能达到最大价值的动作。

也可以写成不确定性集合下的乐观优化形式。定义
\[
\mathcal U_t(a)
=
\left[
\mu_a-\beta_t\sigma_a,\,
\mu_a+\beta_t\sigma_a
\right].
\]
则
\[
\sup_{q\in\mathcal U_t(a)}q
=
\mu_a+\beta_t\sigma_a.
\]

所以 UCB 等价于
\[
a_t^{\mathrm{UCB}}
=
\arg\max_{a\in\mathcal A}
\sup_{q\in\mathcal U_t(a)}
q.
\]

也就是
\[
\boxed{
a_t^{\mathrm{UCB}}
=
\arg\max_{a}
\sup_{q\in\mathcal U_t(a)}
q
}
\]

这说明 UCB 的本质是：

\[
\boxed{
\text{在不确定性约束下，对每个动作取最乐观可能值，然后最大化。}
}
\]

\section{KG：采样后对最大值认知更新的期望}

Knowledge Gradient，简称 KG，不是直接看哪个动作当前上界最大，而是看：

\[
\boxed{
\text{采样哪个动作，最能提升我们对全局最优值的认知。}
}
\]

假设我们现在采样动作 $a$，得到随机观测
\[
Y_a.
\]

观测 $Y_a$ 后，后验均值会更新为
\[
\mu_{t+1}(b;Y_a),
\]
其中 $b\in\mathcal A$ 表示任意候选动作。

更新后的认知最优值为
\[
V_{t+1}(Y_a)
=
\max_{b\in\mathcal A}
\mu_{t+1}(b;Y_a).
\]

当前认知最优值为
\[
V_t
=
\max_{b\in\mathcal A}
\mu_t(b).
\]

因此，采样动作 $a$ 带来的认知提升为
\[
\Delta_a(Y_a)
=
V_{t+1}(Y_a)-V_t.
\]

KG 定义为这个认知提升的期望：
\[
\mathrm{KG}_t(a)
=
\mathbb E_{Y_a}
\left[
V_{t+1}(Y_a)-V_t
\right].
\]

代入 $V_{t+1}$ 和 $V_t$，可得
\[
\mathrm{KG}_t(a)
=
\mathbb E_{Y_a}
\left[
\max_{b\in\mathcal A}
\mu_{t+1}(b;Y_a)
\right]
-
\max_{b\in\mathcal A}
\mu_t(b).
\]

KG 策略为
\[
a_t^{\mathrm{KG}}
=
\arg\max_{a\in\mathcal A}
\mathrm{KG}_t(a).
\]

即
\[
\boxed{
a_t^{\mathrm{KG}}
=
\arg\max_a
\mathbb E_{Y_a}
\left[
\max_b \mu_{t+1}(b;Y_a)
-
\max_b \mu_t(b)
\right].
}
\]

\section{阈值型 KG：最大值认知更新超过某个阈值的概率}

如果我们不只关心平均提升，而是关心：

\[
\text{采样后，对最优值的认知提升是否超过某个阈值 } \tau,
\]

则可以定义阈值型 KG。

仍然令
\[
\Delta_a(Y_a)
=
V_{t+1}(Y_a)-V_t.
\]

那么可以定义概率型指标：
\[
\mathrm{PI\text{-}KG}_t(a)
=
\mathbb P_{Y_a}
\left(
\Delta_a(Y_a)>\tau
\right).
\]

也就是
\[
\mathrm{PI\text{-}KG}_t(a)
=
\mathbb P_{Y_a}
\left[
\max_b \mu_{t+1}(b;Y_a)
-
\max_b \mu_t(b)
>
\tau
\right].
\]

对应的选择策略为
\[
a_t
=
\arg\max_a
\mathbb P_{Y_a}
\left[
\max_b \mu_{t+1}(b;Y_a)
-
\max_b \mu_t(b)
>
\tau
\right].
\]

这个形式可以理解为：

\[
\boxed{
\text{选择那个最有可能让我们对最大值的认知提升超过阈值的动作。}
}
\]

也可以定义风险敏感版本：
\[
\mathrm{Threshold\text{-}KG}_t(a)
=
\mathbb E_{Y_a}
\left[
\left(
V_{t+1}(Y_a)-V_t-\tau
\right)_+
\right],
\]
其中
\[
(x)_+ = \max(x,0).
\]

即
\[
\mathrm{Threshold\text{-}KG}_t(a)
=
\mathbb E_{Y_a}
\left[
\left(
\max_b \mu_{t+1}(b;Y_a)
-
\max_b \mu_t(b)
-\tau
\right)_+
\right].
\]

\section{UCB 与 KG 的核心区别}

\begin{table}[h]
\centering
\begin{tabular}{lll}
\toprule
方法 & 关注的问题 & 数学形式 \\
\midrule
UCB 
& 谁在乐观情况下可能最大 
& $\arg\max_a \mu_a+\beta_t\sigma_a$ \\

KG 
& 采样谁最能提升对最优值的认知 
& $\arg\max_a \mathbb E[V_{t+1}-V_t]$ \\

阈值 KG 
& 谁最可能带来超过阈值的认知跃迁 
& $\arg\max_a \mathbb P(V_{t+1}-V_t>\tau)$ \\
\bottomrule
\end{tabular}
\end{table}

因此，二者的区别可以概括为：

\[
\boxed{
\text{UCB 选择“可能价值最高”的动作。}
}
\]

\[
\boxed{
\text{KG 选择“最能帮助我们知道谁价值最高”的动作。}
}
\]

\section{一个二选项直觉例子}

假设有两个动作 $A$ 和 $B$：
\[
A:\quad \mu_A=10,\quad \sigma_A=1,
\]
\[
B:\quad \mu_B=8,\quad \sigma_B=5.
\]

当前认知下，
\[
V_t=\max(10,8)=10.
\]

若取 $\beta=2$，则
\[
\mathrm{UCB}(A)=10+2\times 1=12,
\]
\[
\mathrm{UCB}(B)=8+2\times 5=18.
\]

因此 UCB 会选择 $B$，因为虽然 $B$ 当前均值较低，但它的不确定性更大，乐观上界更高。

而 KG 会问：

\[
\text{如果我采样 }B\text{，是否会改变我对最优动作的判断？}
\]

如果采样 $B$ 后，其后验均值可能更新到 $11$，那么
\[
V_{t+1}=11,
\]
认知提升为
\[
\Delta_B=11-10=1.
\]

如果采样 $B$ 后，其后验均值只更新到 $8.5$，那么
\[
V_{t+1}=10,
\]
认知提升为
\[
\Delta_B=0.
\]

所以 KG 实际上衡量的是
\[
\mathrm{KG}(B)
\approx
\mathbb E
\left[
(\mu_B^+-\mu_A)_+
\right],
\]
其中 $\mu_B^+$ 表示采样 $B$ 后更新得到的后验均值。

阈值型 KG 则看
\[
\mathbb P
\left(
\mu_B^+-\mu_A>\tau
\right).
\]

\section{最终总结}

UCB 的概率描述是：
\[
\boxed{
a_t^{\mathrm{UCB}}
=
\arg\max_a
\sup_{q\in\mathcal U_t(a)}
q
}
\]
即在不确定性集合中取最乐观值，然后最大化。

KG 的概率描述是：
\[
\boxed{
a_t^{\mathrm{KG}}
=
\arg\max_a
\mathbb E_{Y_a}
\left[
\max_b \mu_{t+1}(b;Y_a)
-
\max_b \mu_t(b)
\right]
}
\]
即选择那个采样后最能提升我们对最优值认知的动作。

阈值型 KG 可以写成：
\[
\boxed{
a_t
=
\arg\max_a
\mathbb P_{Y_a}
\left[
\max_b \mu_{t+1}(b;Y_a)
-
\max_b \mu_t(b)
>
\tau
\right].
}
\]

这正对应：

\[
\boxed{
\text{选择那个最可能让“对最大值的理解更新”超过某个阈值的选项。}
}
\]

\end{document}